\section{AR-19. A direct Katz nonvanishing calculation at
5}\label{ar-19.-a-direct-katz-nonvanishing-calculation-at-5}

Run \texttt{JAC1964-20261121}; \texttt{/root/arithmetic\_sign}, GPT-6
Astra Pro; effort setting not exposed. This is a new deduction from the
explicitly cited Katz-measure construction. The complete proof and exact
script have passed independent whole-read checks by the root and h9
agents. It establishes a p-adic analytic input, not injectivity of a
Chow Abel--Jacobi map.

\subsection{Statement and conventions}\label{statement-and-conventions}

Put \(K=\mathbf Q(i)\), \(\mathfrak f=(4)\), and let \(\psi\) be the
ideal-valued Grossencharacter of \[
E:\quad y^2=x^3+x,\qquad \omega=\frac{dx}{2y}.
\] Thus \(\psi\) is the character denoted \(\psi'\) in BDP §5.2, its
conductor is \((4)\), and \(\psi\bar\psi=\mathbf N\). At the split
ordinary prime \(p=5\), define the two ideal-valued characters \[
\rho=\frac{\bar\psi^2}{\psi},\qquad
\rho_0=\frac{\psi^3}{\bar\psi^2}.
\tag{AR19.1}
\] We use the Bannai--Kobayashi (BK) convention: their Katz value
\(L_{p,\mathfrak f}(\rho_0)\) interpolates \(L(\rho_0^{-1},0)\), and
\(\rho_0\) has infinity type \((3,-2)\) in their convention. Castella's
infinity-type convention is opposite for the same ideal-valued
character, and his Katz function interpolates \(L(\xi,0)\). Thus the
relevant comparison is \(\xi=\rho^{-1}\). In particular \[
\rho^{-1}=\psi^3\mathbf N^{-2}=\chi\mathbf N^{-2},
\] where \(\chi\) is the weight-four Fermat character of AR-11.

\textbf{Proposition AR-19.} With compatible good-reduction CM periods,
the BK value \(L_{p,(4)}(\rho)\) is a 5-adic unit. Consequently the Katz
nonvanishing premise (AR-KATZ) holds: \[
\mathcal L_v(\chi\mathbf N^{-2})(0)\ne0.
\tag{AR19.2}
\] The statement concerns the Katz function itself and does not use the
ramified Rankin--Katz comparison factor in BDP §5.2.

\subsection{1. The exact fourth-order theta
coefficient}\label{the-exact-fourth-order-theta-coefficient}

Write \(E(\mathbf C)=\mathbf C/\Gamma\), \(\Gamma=4\Omega\mathbf Z[i]\),
using the differential \(\omega\), and put
\(A=\operatorname{area}(\mathbf C/\Gamma)/\pi\). The modified weight-two
Eisenstein constant vanishes for a square lattice: multiplication by
\(i\) fixes the lattice and multiplies that constant by \(i^{-2}=-1\).
Hence the reduced theta function is the Weierstrass sigma function with
derivative one at zero.

For a nonzero torsion point \(P=z_P\bmod\Gamma\), put \[
F(P)=\zeta(z_P)-\frac{\bar z_P}{A},\qquad
x=\wp(z_P),\qquad y=\frac{\wp'(z_P)}2.
\] The translated Kronecker function in BK equation (15), p.15, has the
expansion \[
\Theta_{P,0}(z,w)=\frac1w+F(z+P)
 +\frac{F(z+P)^2-\wp(z+P)}2w
 +\frac{F(z+P)^3-3F(z+P)\wp(z+P)-\wp'(z+P)}6w^2+O(w^3),
\] where in this formula \(F(z+P)=\zeta(z+z_P)-\bar z_P/A\) is
holomorphic in \(z\); the conjugate of the variable \(z\) is not
subtracted. Its derivatives satisfy \[
F'=-x,\qquad x'=2y,\qquad y'=3x^2+1.
\] By BK Theorem 1.17, p.17, comparison of the coefficient of \(z^2w^2\)
gives \[
\mathcal E(P):=\frac{e^*_{2,3}(P,0)}{A^2}
 =-F(P)(2x^2+1)-y\bigl(F(P)^2+x\bigr).
\tag{AR19.3}
\] Indeed applying the derivation \(D(F)=-x\), \(D(x)=2y\),
\(D(y)=3x^2+1\) twice to \((F^3-3Fx-2y)/6\) gives the right side. The
two factorials give
\(\partial_z^2\partial_w^2\Theta_{P,0}(0,0)=2\mathcal E(P)\).

Choose a primitive \(4\)-torsion generator with \(2P=(i,0)\). Its two
conjugates over \(K\) can be written \[
P_\pm=\bigl(i(1\pm\sqrt2),\ (1\pm\sqrt2)(1-i)\bigr).
\tag{AR19.4}
\] The duplication formula verifies \(2P_\pm=(i,0)\); adding
\((-i,0)=i(2P_+)\) to \(P_+\) gives \(P_-\). Thus \(P_-=[1+2i]P_+\). The
reduced zeta function vanishes at a nonzero half-period, so its
duplication formula gives \[
F(P_\pm)=-\frac{3x(P_\pm)^2+1}{4y(P_\pm)}.
\] Substitute (AR19.4) into (AR19.3). Using only \(i^2=-1\) and
\((\sqrt2)^2=2\), one obtains the exact identities \[
\mathcal E(P_+)=\mathcal E(P_-)=-\frac{1+i}{4},
\qquad
\mathcal E(P_+)+\mathcal E(P_-)=-\frac{1+i}{2}.
\tag{AR19.5}
\] Changing a lattice generator by a Gaussian unit multiplies this sum
by a Gaussian unit and cannot affect the ensuing nonvanishing. Among the
two points in (AR19.4), at least one is represented by a lattice
generator divided by four: the primitive classes of
\(\mathbf Z[i]/(4)\), modulo global Gaussian units, are precisely its
two ray classes. Both points give the same coefficient, so that choice
is immaterial.

\subsection{2. The two ray classes and their measure
normalizations}\label{the-two-ray-classes-and-their-measure-normalizations}

The ray class group modulo \((4)\) has order two:
\(|(\mathbf Z[i]/4)^\times|=8\) and the four global units inject. Use
inverse ideal representatives \(1\) and \(\mathfrak a^{-1}\), with \[
\mathfrak a=(\beta),\qquad \beta=1+6i,\qquad N\mathfrak a=37.
\] Both are prime to \(5\mathfrak f\), and \(\beta\equiv1+2i\pmod4\).
For this curve \(\psi(\mathfrak a)=\beta\). This also follows from the
usual primary-generator formula \(\psi=\psi_{y^2=x^3-x}\mu\): \(-\beta\)
is primary and the quadratic character of \(K(\sqrt2)/K\) has value
\(-1\) at norm 37. The resulting isogeny scalar in BK §3.3 is \[
\Lambda(\mathfrak a)=\psi(\mathfrak a)=\beta.
\] The ray class field is \(K(\sqrt2)\). We may take the same equation
of \(E\) on that field; all its conjugates have the same differential
and period lattice. BK §3.3, pp.40--42, then gives \[
\Omega_p^{\sigma_{\mathfrak a}}
 =\frac{\Lambda(\mathfrak a)}{N\mathfrak a}\Omega_p,
\qquad \Omega_p\in W(\overline{\mathbf F}_5)^\times.
\tag{AR19.6}
\] The choices \(\alpha_0=1\in\mathfrak a^{-1}\cap P_K(\mathfrak f)\) in
their Definition 3.9 give translated points \(\Omega\) and
\(\Lambda(\mathfrak a)\Omega\). These are \(P_+,P_-\), up to the same
harmless choice of Gaussian unit.

For any such representative \(\mathfrak b^{-1}\), BK Definitions
3.8--3.9, equations (38)--(39), and the displayed formula following
Definition 3.9, pp.41--42, yield at \(\rho_0\) (with \(k=3,l=2\)) \[
L_{p,\mathfrak f}(\rho_0,\mathfrak b^{-1})
 =\Omega_p\rho_0(\mathfrak b)^{-1}(N\mathfrak b)^2(N\mathfrak f)^2
 \int_{\mathbf Z_5^\times\times\mathbf Z_5^\times}
       x^2y^2\,d\mu_{\Omega,0}(x,y;\mathfrak b^{-1}\Gamma).
\tag{AR19.7}
\] These measures are integral; this is Corollary 2.17, p.30, and
Definition 3.2, p.36, applied to prime-to-5 torsion on the good model.
By the lattice-scaling definition of the measure on p.41 and (AR19.6),
its moment over all of \(\mathbf Z_5^2\) is \[
\int_{\mathbf Z_5^2}x^2y^2\,d\mu_{\Omega,0}(x,y;\mathfrak b^{-1}\Gamma)
 =2\Lambda(\mathfrak b)
       (\Omega_p^{\sigma_{\mathfrak b}})^4
       \mathcal E\bigl(\Lambda(\mathfrak b)\Omega\bigr).
\tag{AR19.8}
\] Although the second torsion argument is zero, BK's regularization
subtracts a term depending only on the second formal variable. Two
derivatives in the first variable kill it, so it makes no contribution
to this moment.

The multiplier of \(2\Omega_p^5(N\mathfrak f)^2\mathcal E\) after
substituting (AR19.8) in (AR19.7) is exactly \[
\frac{\Lambda(\mathfrak b)^5}
 {\rho_0(\mathfrak b)(N\mathfrak b)^2}=1,
\tag{AR19.9}
\] since \(\Lambda=\psi\),
\(N\mathfrak b=\psi(\mathfrak b)\bar\psi(\mathfrak b)\), and
\(\rho_0=\psi^3/\bar\psi^2\).

\subsection{3. The unit congruence and the noncritical
value}\label{the-unit-congruence-and-the-noncritical-value}

On the complement of \(\mathbf Z_5^\times\times\mathbf Z_5^\times\) in
\(\mathbf Z_5^2\), the function \(x^2y^2\) is divisible by \(25\).
Integrality of the measure therefore permits replacing the unit-domain
moment in (AR19.7) by (AR19.8) modulo 5. Every exterior factor in
(AR19.7) is a unit. Since \(N\mathfrak f=16\), summing the two classes
and using (AR19.5) gives \[
L_{p,(4)}(\rho_0)
 \equiv2\Omega_p^5\,16^2
          \left(-\frac{1+i}{2}\right)
 =-256(1+i)\Omega_p^5\pmod5.
\tag{AR19.10}
\] This is a unit at either prime over 5: \(1+i\) has norm two, and
\(\Omega_p\) is a unit.

On the Galois group of the ray-class tower of conductor
\(4\cdot5^\infty\), both \(\psi\) and \(\bar\psi\), interpreted as
5-adic characters in the BK convention, take values in
\(\mathbf Z_5^\times\). Therefore \[
\frac{\rho_0}{\rho}
 =\left(\frac{\psi}{\bar\psi}\right)^4
 \equiv1\pmod5.
\tag{AR19.11}
\] Integrality of the Katz measure gives \[
L_{p,(4)}(\rho)\equiv L_{p,(4)}(\rho_0)\not\equiv0\pmod5.
\tag{AR19.12}
\] This proves the BK unit assertion without a numerical approximation.

Finally choose the \textbf{same compatible complex and 5-adic period
data} in BK and Castella, attached to the good equation and differential
fixed above. Compare BK Theorem 3.11, pp.44--46, with Castella Theorem
2.2.1, pp.7--8. Under \(\xi=\rho^{-1}\), their interpolation values
involve the same Hecke value, gamma factor, and the two Euler factors at
5. BK writes one factor explicitly and the other in
\(L_{\mathfrak f\bar{\mathfrak p}}\). In particular \(\rho_0^{-1}\) has
Castella infinity type \((3,-2)\), so that the displayed formulas match
literally with this period choice. More generally they match on the
critical characters in the component with \(k\equiv3\pmod4\),
\(l\equiv2\pmod4\); the BK sign \((-1)^{k+l-1}\) is then +1. These
critical characters are dense in that two-variable component. The
uniqueness of the integral Katz measure therefore identifies the two
functions there with this inverse-character convention. Both \(\rho\)
and \(\rho_0\) lie in the component by (AR19.11). This proves (AR19.2)
from (AR19.12). Changing the good integral period data multiplies the
evaluation at this fixed algebraic weight by a nonzero unit, so
nonvanishing is independent of that permitted choice; no assertion about
an arbitrary weight-dependent global analytic rescaling is needed.

\subsection{Scope}\label{scope}

AR-19 supplies the previously missing nonzero Katz value for the Fermat
character at 5. Together with a separately checked Euler-system and
duality argument at the actual conductor, it can supply a
Selmer-dimension bound. It does not give scalar-extended étale
Abel--Jacobi injectivity, does not identify a Selmer dimension with a
Chow dimension, and does not prove real-height factorization through a
p-adic Abel--Jacobi map.

\subsection{Primary locators and computational
check}\label{primary-locators-and-computational-check}

\begin{itemize}
\tightlist
\item
  Bannai--Kobayashi, \emph{Algebraic theta functions and the p-adic
  interpolation of Eisenstein--Kronecker numbers}, arXiv:math/0610163v4
  (11 December 2007), equation (15) p.15; Theorem 1.17 p.17; Corollary
  2.17 p.30; Proposition 3.1 and Definition 3.2 pp.35--36; §3.3
  pp.40--43, Definitions 3.8--3.9; Theorem 3.11 pp.44--46. Primary PDF
  directly read, including the period and Euler-factor formulas. Local
  source \texttt{work/h9/BK2007.pdf/.txt}.
\item
  Castella, \emph{Tamagawa number conjecture for CM modular forms and
  Rankin--Selberg convolutions}, 14 February 2025 manuscript, §2.1.1
  pp.5--6 (infinity types and reciprocity), §2.1.2 p.6 (periods),
  Theorem 2.2.1 pp.7--8 (the integral Katz function and uniqueness).
  This use is outside §5 and does not invoke its exact-conductor
  hypothesis.
\item
  BDP, \emph{p-adic L-functions and the coniveau filtration on Chow
  groups}, §5.2, for the character of \(y^2=x^3+x\), its conductor and
  the quadratic twist relation.
\item
  The rational algebra in (AR19.3)--(AR19.5) is also checked by the
  reproducible exact script \texttt{fermat\_katz\_coefficient.py}; no
  floating-point arithmetic is used.
\end{itemize}
